\documentclass[11pt,a4paper]{ctexart}
\usepackage[margin=2.4cm]{geometry}
\usepackage{amsmath,amssymb}
\usepackage{fancyhdr}
\usepackage{enumitem}
\linespread{1.18}
\setlength{\headheight}{15pt}
\setlength{\parskip}{0pt}

\pagestyle{fancy}
\fancyhf{}
\fancyhead[L]{\small Team 2210307}
\fancyhead[R]{\small \thepage}
\renewcommand{\headrulewidth}{0.4pt}
\fancypagestyle{plain}{\fancyhf{}\renewcommand{\headrulewidth}{0pt}}

\ctexset{
  section/format = \Large\bfseries\raggedright,
  subsection/format = \large\bfseries\raggedright,
  subsubsection/format = \normalsize\bfseries\raggedright
}
\setcounter{secnumdepth}{3}

\begin{document}

%====================== 摘要页 ======================
\thispagestyle{empty}
\begin{center}
{\small 2022 MCM/ICM\quad 选题 A\quad 队伍控制号 2210307\quad 摘要页（Summary Sheet）}\\[18pt]
{\LARGE\bfseries 功率规划模型：骑手的制胜法宝}\\[8pt]
{\large （Power Planning Model: Magic Weapon for Cyclists）}\\[10pt]
\end{center}

\begin{center}
{\large\bfseries 摘\quad 要}
\end{center}

\vspace{6pt}
在计时赛中，许多因素都会影响骑手的表现。为帮助骑手取得更好的成绩，我们需要建立一个能够综合考虑众多因素、并在赛道不同位置给出合理功率分配方案的模型。

首先，我们建立了骑手能力评价模型（Rider Ability Evaluation Model），并定义了骑手的功率剖面。根据功率与该功率水平下最长持续时间之间的关系，绘制出功率—时间散点图；随后依据功率—时间关系划分出四个不同的运动强度域。我们选用三参数临界功率模型（3P-CP 模型）对散点图进行拟合，得到骑手的功率—时间曲线，进而得到三项重要的骑手能力指标：最大瞬时功率（$P_{\max}$）、临界功率（$CP$）与无氧能量储备（$W'$）。我们取多组数据对拟合结果进行评价，经计算，拟合度稳定在 0.86 左右，表明拟合效果良好。

其次，我们建立了外部因素分析模型（External Factor Analysis Model），以识别外部条件的潜在影响。我们从路面滚动阻力与风阻两个方面，对天气和环境的影响进行定量评价。此后，我们选取实际赛道进行验证，符合程度超过 90\%。

第三，建立了功率总体规划模型（Power Overall Planning Model）。在对赛道进行分段与动力学分析之后，提出平均功率（AP）与归一化功率（NP），并获得各路段功率的约束条件，从而将问题转化为以总时间 $t$ 为目标函数的多元非线性规划问题。我们采用梯度下降法逐步逼近 $t$ 的局部极小值，并给出各路段功率的分配方案。我们的结果与实际值的偏差为 $\pm 10\%$，说明与实际吻合良好。

随后，我们将骑手能力评价模型应用于不同类型的骑手与赛道，得到骑手的输出功率、特定位置的速度以及完成比赛的最短时间。

之后，我们对模型进行了敏感性分析，分析其对天气与环境以及骑手表现的稳健性。气温、气压、海拔与风阻的失效概率分别为 0.18\%、3.97\%、0.18\% 与 4.22\%。气压与风阻属于敏感因素，而温度与海拔属于非敏感因素。我们计算了偏离功率计划所引起的失效概率，并给出策略执行的优先次序：坡道 $>$ 弯道 $>$ 平路。

最后，我们将模型推广为团队计时赛功率分配模型，得到扩展后的多元非线性规划问题。我们的建议是：赛道前 45\% 由两名爆发力强的骑手领骑，最后 15\% 由耐力最强的骑手领骑，其余 4 名骑手在中段轮流领骑。

\vspace{6pt}
\noindent\textbf{关键词：}计时赛；功率总体规划；多元非线性规划

\newpage
\pagestyle{fancy}

%%====================== 正文 ======================
\section{引言}

\subsection{背景}

“更高、更快、更强——更团结”是奥林匹克格言，也是每一位运动员的目标。为展现拼搏进取的精神，运动员努力取得更好的成绩。然而，影响骑手表现的因素众多，包括赛事类型、赛道以及骑手的能力。为帮助这些骑手应对不同的情形并取得更好的成绩，我们需要制定合理的方案。

\subsection{问题重述}

\begin{enumerate}
\item 为确定骑手在赛道上的位置与其所施加功率之间的关系，我们建立一个可适用于任意类型骑手与任意赛道的模型。
\item 定义两种不同性别骑手的功率剖面。
\item 将我们的模型应用于多种计时赛赛道，并自行设计一条计时赛赛道。
\item 确定我们的模型对天气与环境以及对偏离目标功率的敏感性。
\item 扩展我们的模型，将其用于一支团队。
\end{enumerate}

\subsection{我们的工作}

\begin{itemize}
\item 对于问题 1），我们拟提出一个模型来度量不同类型骑手之间的差异，对赛道进行分段并将其概括为若干特征。我们可以通过用离散点拟合曲线的方式计算参数。我们需要对不同赛道进行动力学分析。对于功率分配的总体规划，可以尝试使用多元非线性规划方法。最终，在已知骑手、环境、赛道等信息之后，得到功率分配的最优方案。
\item 对于问题 2），即模型在不同类型骑手上的应用，我们度量不同骑手在各方面的能力值，并将其代入模型。
\item 对于问题 3），即模型在不同赛道上的应用，我们需要控制变量，使用同一名运动员与三组已知的赛道信息，推导出三种不同的功率分配。
\item 对于问题 4）与问题 5），这两个问题都属于敏感性分析。我们计算不同变量的失效概率，并基于用子集模拟重要性采样的思想度量敏感性。
\item 对于问题 6），我们需要找出团队赛与个人赛之间的异同，并讨论如何给出最优的团队策略。
\end{itemize}

\section{假设与论证}

为了更充分地利用关于骑手与赛道的数据，并提高模型的有效性，我们提出如下合理假设：

\begin{itemize}
\item \textbf{安全第一。}运动员的身体健康与安全是一切的前提。我们假设骑手存在速度、能量与体力上限，以避免模型将他们置于危险的过度透支状态。
\item \textbf{次要因素的影响可忽略不计。}由于比赛成绩受众多因素影响，而我们的模型主要关注不同类型的骑手，因此有必要忽略来自次要因素（如骑手的装备）的影响。
\item \textbf{赛道分段。}我们将整条赛道划分为 $n$ 个路段，每个路段属于以下四种类型之一：平路、上坡、下坡与急弯。分段的判定标准为：当坡度的连续变化不超过 1\% 时，定义为同一段平地；当曲率半径超过 40\,m 时，视为直道；当曲率半径小于 40\,m 且变化不超过 5\,m 时，定义为同一个弯道。
\item \textbf{体育精神。}所有参赛者都具备体育精神，维护比赛的公平性，并能在比赛中发挥出最佳水平。
\end{itemize}

\section{符号说明}

\begin{itemize}
\item $e$：自然对数的底，$e=2.178$；
\item $R$：气体常数，$8.31\ \mathrm{J\cdot mol^{-1}\cdot K^{-1}}$；
\item $g$：重力加速度，$g=9.8\ \mathrm{m/s^{2}}$；
\item $CP$：临界功率（critical power），单位瓦特（watt）；
\item $MMP$：平均最大功率输出（mean maximal power output），单位瓦特；
\item $W'$：无氧能力（anaerobic capacity），单位焦耳（J）；
\item $\rho$：空气密度，单位 $\mathrm{g/L}$；
\item $k_r$：滚动阻力系数；
\item $c_d$：风阻系数；
\item $S$：迎风面积，单位 $\mathrm{m^{2}}$。
\end{itemize}

\section{模型准备}

为使本文更清晰、更易于理解，我们在本节中对一些名词、数据类型与数据来源加以说明。

\subsection{名词解释}

\begin{itemize}
\item \textbf{归一化功率输出（Normalized power output，NP）}

它考虑了稳态训练与波动训练之间的差异。NP 能更好地量化可变用力下“更吃力”感受的生理“代价”。对于波动很大的训练，NP 可远高于平均功率（AP）；而对于非常平稳的训练，NP 与 AP 几乎相等。相对较高的 NP 表明该运动的波动较大，在生理上比平均功率所代表的更为艰难。

\item \textbf{平均功率输出（Average power output，AP）}

骑手的平均功率输出是其所消耗的总能量与完成长距离比赛所用时间之比。AP 的计算方法如下：
\[
AP=\frac{\text{总能量}}{\text{总时间}}
\]

\item \textbf{平均最大功率输出（Mean maximal power output，MMP）}

MMP 是骑手在给定时间内能够产生的功率量。例如，在 5 分钟时长内记录到的最高平均功率输出即为 5 分钟 MMP。此类 MMP 数据很有价值，因为它们可以识别骑手所需的能量输出与持续时间。

\item \textbf{无氧能力（Anaerobic capacity，$W'$）}

$W'$ 表示临界功率以上的做功能力。

\item \textbf{临界功率（Critical power，CP）}

CP 是功率曲线的下限（或称水平渐近线）。它是一项生理阈值，划分了运动中可持续代谢的用力与不可持续的用力。临界功率曲线与 $W'$ 的关系为：红色曲线即临界功率曲线。

\item \textbf{最大瞬时输出功率（Maximum instantaneous output power，$P_{\max}$）}

$P_{\max}$ 是 1 秒内的峰值功率，反映骑手的爆发力。

\item \textbf{最大摄氧量（$VO_{2\max}$）}

$VO_{2\max}$ 指机体在最大强度运动时能够摄入的氧气量。由于氧气是体内众多物质代谢与能量代谢活动的基础，$VO_{2\max}$ 是衡量骑手耐力（stamina）的一项重要指标，反映机体的有氧运动能力。

在该能量消耗过程中，骑手会依次到达乳酸阈、临界功率与最大摄氧量；也就是说，随着能量消耗的推进，若干生理阈值会被逐一触及，而这正是划分运动强度域的依据。
\end{itemize}

\subsection{数据准备}

\begin{itemize}
\item \textbf{赛道条件}

我们从网站\ [1][2] 找到了所需的全部数据，以及相应的天气与环境条件。

\item \textbf{骑手类型}

按题目所给出，骑手被分为五种类型，即计时赛专家（time trial specialist）、爬坡手（climber）、冲刺手（sprinter）、全能手（rouleur）与爆发型选手（puncheur）。我们对这五类骑手作简要描述并归纳其特征，具体如下。

\begin{itemize}
\item \textbf{计时赛专家}：经验丰富、能力出众的计时赛车手。乳酸阈较高，肌肉耐力出色，通常 FTP 值较大，擅长最大化地利用体能（体能评估与分配），并且耐力极强。
\item \textbf{爬坡手}：专门擅长长距离爬坡赛的骑手。体重轻，耐力与腿部力量出色。
\item \textbf{冲刺手}：专门擅长在短时间内产生极限功率的骑手。爆发力强，耐力较差。
\item \textbf{全能手}：可适应多种地形的全能型骑手。各项指标相对均衡，但均非顶尖。
\item \textbf{爆发型选手}：专门擅长包含大量短陡坡或急剧加速的比赛的骑手。体重大，腿部力量强，耐力强，爆发力中等，腰部力量较弱。
\end{itemize}
\end{itemize}

骑手产生的功率越大，其能够维持该功率的时间就越短，之后必须进行恢复。我们用 $t_{maintain}-P$ 图来描述骑手在这一身体条件上的差异，其中 $P$ 为作用于踏板的持续功率水平，$t_{maintain}$ 是骑手在维持当前功率水平的情况下耗尽全部能量所需的时间。

\begin{itemize}
\item \textbf{功率曲线}

要求骑手完成 3 至 5 次不同时长的全力输出，并记录骑手的功率。每次全力输出之间的恢复时间应设置为至少 1 小时，以确保骑手获得足够的恢复时间。作为一个示例，我们得到了骑手在不同全力输出时长下 MMP 的散点图。其目的是对临界功率（CP）等指标进行拟合与回归分析。
\end{itemize}

\section{功率总体规划模型}

\subsection{模型概述}

骑手能力以及其他外部因素（包括道路条件、天气条件等）构成模型的输入参数，骑手在各路段的功率输出则作为输出参数。该模型帮助骑手以最短的时间完成赛道骑行，这需要两个步骤。步骤一：评价骑手的能力，使用我们的骑手能力评价模型。步骤二：对赛道进行分段，对不同类型的路段进行分类并逐一分析，随后得到影响骑手最终表现的各路段功率与时间之间的关系。

\subsection{子模型一：骑手能力评价模型}

\subsubsection{子模型一概述}

我们以功率作为衡量骑手能力的首要指标。通过对不同时长的多次测试，可以绘制出描述骑手最大持续功率与最长持续时间之间关系的散点图，我们认为它可以拟合为一条光滑曲线。拟合曲线为何可以是光滑的？由于我们体内的生理响应是一系列连续的过程，功率—持续时间曲线本身便是光滑的。由于拟合曲线中只有一部分与实际情形吻合良好，我们进行运动强度域划分，以清晰地描述曲线的哪一部分更接近实际情形。

\subsubsection{运动强度域划分}

根据 Burnley 与 Jones 在 2007 年的研究\ [4]，从生理学的角度看，功率—持续时间关系由四个不同的运动强度域构成，即中等（moderate）、较重（heavy）、重度（severe）与极限（extreme）。作出这一划分的依据是：对每一个运动强度域而言，其功率—持续时间曲线都存在相对明显的差异。在骑手全程功率输出图中，我们依据功率输出划分出这四个运动强度域，并用不同的颜色加以表示。

\subsubsection{三参数临界功率模型}

我们选择三参数临界功率模型作为骑手能力评价模型。在 Peter Leo 的一篇论文\ [3] 中，比较了五种描述功率—时间关系的模型，其中 3-P CP 模型（三参数临界功率模型）在重度与极限运动强度域中的结果与实际情形最为吻合。因此，在这两个域中的功率输出对于预测骑手表现更为重要，于是我们决定采用 3-P CP 模型作为功率—时间关系模型来拟合测试数据。由拟合曲线，我们可以得到骑手的功率曲线以及骑手能力指标。

骑手的生理与能力指标如下：身高 $h$ 为骑手的站立身高，单位为米（m）；体重 $m$ 的单位为千克（kg）；爆发力以 $P_{\max}$ 表示，采用骑手在给定 5\,s 内能够输出的平均最大功率（MMP）来衡量骑手的爆发力，单位为瓦特；耐力以 $W$ 表示，我们用无氧能力衡量骑手的耐力，即无氧条件下储存的能量；摄氧能力以 $CP$ 表示，我们用临界功率（CP）表示骑手的摄氧能力，CP 代表功率的理论渐近线，表明给定的功率输出可以一直维持；控车能力以 $BC$ 表示，控车能力越强，骑手的最高速度与攻击性上限就越高。

不足 2 分钟的预测性测试无法保证达到 $VO_{2\max}$（即它们不属于重度强度范围）。近期的工作提出，CP 能更好地估计最大代谢稳态，即在血乳酸值升高的情况下，仍可观察到摄氧量（$VO_{2}$）响应处于稳态时的最高功率输出。

\textbf{功率—时间关系模型。}功率与时间的关系如下式所示：
\[
t=\frac{W'}{P-CP}+\frac{W'}{CP-P_{\max}}
\]
其中 $t$ 为持续时间；$P$ 为最大持续功率；$W'$ 为无氧能量；$CP$ 为临界功率；$P_{\max}$ 为最大瞬时功率。

该模型的拟合通过 MATLAB 拟合工具箱中自定义参数函数的拟合工具实现，拟合结果将在下一部分给出。

在实际中，如果实际条件不足以获得更多的功率—时间散点值，可以通过简化实验测量功率—时间曲线的三个参数，从而同样得到功率—时间曲线。CP 可以通过测量 3 分钟全力输出功率曲线最后 30\,s 的 AP 得到，$P_{\max}$ 可以通过测量 5\,s 全力输出功率的平均值得到。在测得 CP 之后，通过特定输出功率时长的测试即可得到 MMP 与 $t$。此时 $W'$ 可由下式表示：
\[
W'=(MMP-CP)\times t
\]
值得说明的是，尽管上述方法更为简便，但由此测得的功率—时间曲线会比拟合结果存在更大的误差。

\subsubsection{拟合结果评价}

为评价拟合效果，我们对由某运动员多次测试数据绘制的散点图进行拟合，拟合结果如下所述。其拟合优度（Goodness of Fit）为 0.864，接近理想值 1，因此拟合良好，模型具有实际意义。

在曲线拟合之后，我们可以得到该光滑曲线的表达式，表达式中的每个参数即为骑手的能力指标。

\subsection{子模型二：外部因素分析模型}

\subsubsection{天气与环境因素}

天气与环境对比赛的影响是复杂的，主要体现在路面因素与空气阻力因素两个方面。骑手所受的阻力主要来自地面滚动阻力与风阻，骑手往往需要付出全部功率的 40\% 来克服这两种阻力。因此，我们从道路条件与风阻两个方面来考虑天气与环境因素的影响。

\begin{itemize}
\item \textbf{道路条件}

滚动阻力主要由地面的弹力与滚动阻力系数决定，其表达式为：
\[
F_r=k_r\times F
\]
其中 $F_r$ 为滚动阻力，$F$ 为地面的弹力，$k_r$ 为滚动阻力系数，它与路面类型、轮胎结构等有关。

\item \textbf{风阻}

\begin{itemize}
\item \emph{风的强度与风向。}风阻是高速骑行时风力在骑行方向上的分量，主要与风速和风向有关。

无风时，风阻的计算公式为：
\[
F_w=\frac{1}{2}\times c_d\times S\times \rho\times v^{2}
\]
其中 $v$ 为骑手的速度，$c_d$ 为风阻系数，$S$ 为迎风面积。

有风时，假设在短时间内风速与风向固定，风向与骑行方向的夹角为 $\theta$。以风为参考系时，骑手处于无风状态，骑手的速度 $v_r$ 为：
\[
v_r=v-v_w\times\cos\theta
\]
其中 $v_w$ 为风速。于是风阻为：
\[
F_w=\frac{1}{2}\times c_d\times S\times \rho\times v_r^{2}
\]
考虑风的方向，则当 $0^{\circ}<\theta<90^{\circ}$ 时风阻 $>0$，当 $90^{\circ}<\theta<180^{\circ}$ 时风阻 $<0$，于是最终的风阻表达式为：
\[
F_w=\frac{1}{2}\times c_d\times S\times \rho\times \frac{v_r^{3}}{\left|v_r\right|}
\]

\item \emph{空气密度。}空气密度 $\rho$ 与温度和气压有关，其计算公式为：
\[
\rho=\frac{pM}{RT}
\]
其中 $p$ 为气压，$M$ 为气体的摩尔质量，$R$ 为比例常数（对任意理想气体均为定值），$T$ 为气温。由于气压往往不易直接测量，我们将气压的变化转换为海拔的变化，二者的关系如下：
\[
H_s=H_0+\frac{R}{g}\times T_m\times \ln\frac{P_0}{P_s}
\]
其中 $H_0$ 为海拔高度，$H_s$ 为标准等压面的海拔高度，$P_0$ 为地面气压，$P_s$ 为气柱的平均高度，$R$ 与 $g$ 为常数。

\item \emph{风阻系数。}阻力系数仅由骑手的骑行姿态决定。基于前述假设，我们认为该参数保持不变。对于自行车骑手，我们取 $c_d=0.23$。

\item \emph{迎风面积。}根据 Stevenson 公式，迎风面积与身高、体重有关，其表达式为：
\[
S_{\max}=
\begin{cases}
\dfrac{1}{4}(0.0057h+0.0121m+0.0882), & \text{男骑手}\\[0.5em]
\dfrac{1}{4}(0.0073h+0.0127m-0.0099), & \text{女骑手}
\end{cases}
\]
\end{itemize}
\end{itemize}

\subsubsection{计时赛赛道因素}

我们的模型还需要适用于任意类型的赛道。因此，为使赛道能够被定量分析，我们需要归纳其特性。于是我们将赛道的所有路段划分为四种类型：弯道、上坡、下坡与平路。
\subsection{建立功率总体规划模型}

\subsubsection{赛道上的动力学分析}

我们用参数来度量四种类型的赛道路段：平路（参数为距离、速度与功率）、上坡（参数为距离、坡度、速度与功率）、下坡（参数为距离、坡度、速度与功率）、弯道（参数为距离、曲率半径、速度与功率）。下面分别加以讨论。

\begin{itemize}
\item \textbf{平路}

\[
P=(F_w+F_r)\times v=mg\times k_r\times v+\frac{1}{2}\times c_d\times S\times \rho\times v^{3}
\]
其中 $F_w$ 为风阻（即受力图中的空气阻力），$F_r$ 为滚动阻力。

\item \textbf{上坡}

\[
P=F\times v=(F_r+F_g+F_w)\times v=\left(k_r\times mg\times \cos\theta+mg\times \sin\theta\right)v+\frac{1}{2}\times c_d\times S\times \rho\times v^{3}
\]
其中 $F_g$ 为骑手的重力（即受力图中的 gravity），且角度 $\theta>0$。

\item \textbf{下坡}

我们将坡度超过 2\% 的路段定义为下坡。由于下坡路段与上坡路段具有相似性，上坡时 $\theta>0$，下坡时 $\theta<0$，因此这里只给出上坡的受力公式，下坡的受力公式与之类似。

\item \textbf{弯道}

\[
\begin{cases}
P=\left(F_w+\sqrt{F_r^{2}-F_c^{2}}\right)\times v\\[0.5em]
F_c=\dfrac{m\times v^{2}}{r}
\end{cases}
\]
其中 $F_c$ 为向心力。
\end{itemize}

\subsubsection{全程功率输出评价}

在分析骑手在比赛中的表现时，我们使用如下指标来评价骑手在比赛全程的功率输出。

\begin{itemize}
\item \textbf{平均功率（AP）}，计算公式为
\[
AP=\left[\frac{\sum\left(P_i^{4}\times t_i\right)}{t}\right]^{\frac{1}{4}}
\]
平均功率是整段骑行中的总功除以总骑行时间，与速度和力具有很强的相关性。

\item \textbf{归一化功率（NP）}，计算公式为
\[
NP=\frac{\sum P_i\times t_i}{t}
\]
它是对过高与过低的功率点进行平滑处理、并提高高功率输出区间所占比重之后得到的结果，比平均功率更能描述整个运动过程。

\item \textbf{总功（$W_t$）}，计算公式为
\[
W_t=AP\times t
\]
即平均功率与总时间的乘积。它仅反映移动过程中所做的功，不考虑由功率变化引起的额外做功消耗。

\item \textbf{归一化功（$W_n$）}，计算公式为
\[
W_n=NP\times t
\]
归一化功是归一化功率与总时间的乘积，用于反映比赛中实际所做的功，其中包括由功率变化引起的额外做功消耗，通常大于总功。

\item \textbf{最大瞬时功率（$P_{\max}$）}，它反映骑手在比赛中输出功率的上限。该上限不应与平均功率相差过大，否则会因功率变化带来的额外功率消耗而无法达到最优情形。
\end{itemize}

平均功率的缺点在于：如果骑行过程中功率波动较大，算术平均功率不足以反映整个骑行强度。例如，高强度冲刺训练的 AP 过低，无法体现其高强度，而归一化功率 NP 则更好。

如果骑手在比赛中功率输出不断变化，他将消耗额外的功。因此，若非出于心理策略等原因必须不断变化，我们不建议这样做，因为这不利于骑手取得好成绩。

\subsubsection{多元非线性规划}

经过上述分析，我们的问题转化为一个多元非线性规划问题，目标函数为总时间 $t$，且 $t$ 需要被最小化。其数学表达式为：
\[
\min t=\sum_{i=1}^{n}t_i=\sum_{i=1}^{n}\frac{d_i}{v_i}
\]
\[
\text{s.t.}\quad
\begin{cases}
1.15\,MMP(t_i)>P_i>0,\\[0.3em]
NP\leqslant MMP(t),\\[0.3em]
CP\leqslant MMP(t),\\[0.3em]
MMP(t_i)=CP+\dfrac{1}{t}W'-\dfrac{1}{CP-P_{\max}},\\[0.6em]
NP=\dfrac{\sqrt[4]{\sum_{i=1}^{n}P_i^{4}t_i}}{t},\\[0.6em]
AP=\dfrac{\sum_{i=1}^{n}P_it_i}{t},\\[0.6em]
P_i\geqslant 0,\ t_i>0.
\end{cases}
\]
该问题是一个典型的多元非线性规划问题。我们求解各路段的最优输出功率 $P_i$，以实现总时间 $\min(t)$ 最小。

\subsection{模型求解与评价}

\subsubsection{模型求解}

由于多元非线性规划问题难以求得精确的最小解，我们考虑使用梯度下降法逐步逼近局部最优解以寻找最优解。在比较多个局部最优解之后，可以得到一个近似全局最优解，即较好的解。求解通过 Python 的 optimize 工具箱中的 minimize 函数完成。

值得说明的是，我们无法确定下一个局部最优是否就是我们所寻找的解，因此通常会设定若干次循环，在达到设定次数后停止搜索最优解。当前所有局部极小值中的最小值即为全局最优值。

\subsubsection{模型求解方法评价}

Python 的 optimize 工具箱中的 minimize 函数给出了多种求解多元非线性规划的方法，包括 'SLSQP'、'Nelder-Mead'、'Powell'、'CG'、'BFGS'、'Newton-CG'、'L-BFGS-B' 等。通过比较不同方法，我们发现 'SLSQP' 求解更快，即时间复杂度更低，因此我们选择 'SLSQP' 方法。

\section{模型的应用：骑手与赛道}

\subsection{骑手：不同性别与类型}

我们绘制了四类骑手能力的雷达图：左侧图给出计时赛专家与爆发型选手的六项能力，中间图给出不同性别之间的差异，右侧图给出四类骑手的能力。由雷达图的面积大小可以直观地比较不同性别、不同类型骑手在六项能力指标上的整体强弱。

\subsection{赛道：三条计时赛赛道}

对于不同的计时赛赛道，我们在应用模型时采用控制变量法，即我们的模型只研究同一名骑手在不同赛道上比赛时的最优功率分配策略。同时，我们只考察女子个人计时赛的赛道条件，该骑手的信息为：性别为女性，身高 165\,cm，体重 55\,kg，爆发力 $P_{\max}=990\ \mathrm{W}$，耐力 $W'=550\ \mathrm{J}$，摄氧能力 $CP=230\ \mathrm{W}$，控车能力 0.87。

\subsubsection{2021 年日本东京奥运会计时赛赛道}

尽管距离不同，男子与女子比赛的计时赛赛道完全相同。我们从弯道与坡度两个方面对赛道进行分析与分段。赛道信息为：气温 $15.7\,^{\circ}\mathrm{C}$，气压 0.96\,atm，平均海拔 560\,m，风向为北偏西 $15^{\circ}$，风速 7.5\,m/s，空气湿度 60\%。赛道的环境因素为：滚动阻力系数 $k_r=0.015$，空气密度 $\rho=1.2\ \mathrm{kg/m^{3}}$，风阻系数 $c_d=0.23$。

该赛道的急弯较多，赛道为富士国际赛车场，共有七个急弯，弯道旁的黑色数字代表曲率半径，绿色数字代表弯道长度。赛道高程剖面显示，我们依据不同的坡度将赛道划分为若干路段，并以相同的颜色标注各路段的坡度与长度。

计算所得的赛道数据如下（依次为路段编号、坡度、长度、输出功率、速度与时间）：

\begin{itemize}
\item 路段 1：坡度 $-2.5\%$，长度 2\,km，输出功率 181.34\,W，速度 32.36\,km/h，时间 3.68\,min；
\item 路段 2：坡度 $-5.5\%$，长度 1.7\,km，输出功率 157.35\,W，速度 44.84\,km/h，时间 2.27\,min；
\item 路段 3：坡度 $6.5\%$，长度 0.3\,km，输出功率 200.10\,W，速度 34.42\,km/h，时间 0.52\,min；
\item 路段 4：坡度 $-4.5\%$，长度 0.85\,km，输出功率 161.91\,W，速度 36.20\,km/h，时间 1.41\,min；
\item 路段 5：坡度 $4\%$，长度 5.45\,km，输出功率 193.21\,W，速度 24.48\,km/h，时间 13.3\,min；
\item 路段 6：坡度 $-3.5\%$，长度 4.15\,km，输出功率 158.85\,W，速度 34.19\,km/h，时间 10.29\,min；
\item 路段 7：坡度 $0\%$，长度 1.15\,km，输出功率 183.63\,W，速度 31.95\,km/h，时间 2.15\,min；
\item 路段 8：坡度 $5\%$，长度 1.8\,km，输出功率 223.97\,W，速度 26.49\,km/h，时间 4.07\,min；
\item 路段 9：坡度 $-2.2\%$，长度 2.6\,km，输出功率 189.63\,W，速度 33.28\,km/h，时间 4.68\,min；
\item 路段 10：坡度 $3\%$，长度 2.1\,km，输出功率 190.28\,W，速度 32.77\,km/h，时间 3.84\,min。
\end{itemize}

求解后，最短时间为 46.21\,min，赛道总长 22.1\,km（13.7\,mi），赛道平均坡度 3.5\%，坡度较大，这对骑手的持续输出功率提出了更高要求，因此爬坡手更具优势，因为爬坡手的临界功率（CP）与无氧能量（$W'$）更大。

为了更细致地展示功率分配结果，我们将功率、速度、高程与距离绘制在同一张图上，横轴为距离（mi），三条不同颜色的曲线分别代表功率（watts）、速度（km/h）与高程（ft）。当我们把赛道划分为更短的路段时，也可以得到同样的功率输出、速度、高程与距离对应关系图。

可以看出，为使总时间最小，分配给坡度较小路段的功率应当较低；分配给坡度较大路段的功率应当较高。上坡时使用高输出功率所带来的时间减少量大于下坡时，因此上坡所需的输出功率应大于下坡。

\subsubsection{2021 年比利时佛兰德 UCI 世界锦标赛计时赛赛道}

该赛道信息为：气温 $28\,^{\circ}\mathrm{C}$，气压 1.013\,atm，平均海拔 10\,m，风向为东偏南 $10^{\circ}$，风速 11\,m/s，空气湿度 87\%。赛道的环境因素为：滚动阻力系数 $k_r=0.012$，空气密度 $\rho=1.23\ \mathrm{kg/m^{3}}$，风阻系数 $c_d=0.241$。

该赛道的急弯与坡度分段特征同样由赛道高程剖面给出。计算所得的赛道数据如下：

\begin{itemize}
\item 路段 1：坡度 $0.3\%$，长度 1.65\,km，输出功率 204.58\,W，速度 27.34\,km/h，时间 3.62\,min；
\item 路段 2：坡度 $-0.5\%$，长度 1.55\,km，输出功率 127.93\,W，速度 32.96\,km/h，时间 2.82\,min；
\item 路段 3：坡度 $-1.35\%$，长度 0.8\,km，输出功率 100.04\,W，速度 47.00\,km/h，时间 1.02\,min；
\item 路段 4：坡度 $0\%$，长度 7.15\,km，输出功率 165.97\,W，速度 31.83\,km/h，时间 13.47\,min；
\item 路段 5：坡度 $1\%$，长度 0.85\,km，输出功率 209.77\,W，速度 22.19\,km/h，时间 2.29\,min；
\item 路段 6：坡度 $-0.45\%$，长度 2.00\,km，输出功率 131.45\,W，速度 36.82\,km/h，时间 3.25\,min；
\item 路段 7：坡度 $0\%$，长度 8.10\,km，输出功率 167.90\,W，速度 33.01\,km/h，时间 14.68\,min；
\item 路段 8：坡度 $0.15\%$，长度 7.2\,km，输出功率 184.38\,W，速度 27.65\,km/h，时间 16.26\,min。
\end{itemize}

求解后，最短时间为 57.41\,min，赛道总长 29.3\,km（18.2\,mi），赛道平均坡度 0.17\%。坡度极为平缓，而冲刺手具有更强的爆发力（$P_{\max}$），因此在该赛道上具有优势。由于没有陡上坡，冲刺手可以把全部功率用于相对平坦路面上的快速冲刺，且功率输出分布较为均匀。

所得的功率、速度、高程与距离对应关系同样以图形方式给出，可以直观地看到功率在平坦路段上的均匀分配。

\subsubsection{自行设计的赛道}

我们自行设计的赛道信息为：气温 $15\,^{\circ}\mathrm{C}$，气压 0.91\,atm，平均海拔 345\,m，风向为东偏南 $25^{\circ}$，风速 7.5\,m/s，空气湿度 55\%。赛道的环境因素为：滚动阻力系数 $k_r=0.019$，空气密度 $\rho=1.17\ \mathrm{kg/m^{3}}$，风阻系数 $c_d=0.23$。

我们自行设计的赛道总长 13.78\,km（8.56\,mi），其特点是极陡且起伏崎岖，最大坡度达 25\%，平均坡度 11\%。在这类赛道上，爬坡手很容易在爬坡环段甩开对手。计算得到的最短时间为 24.58\,min。该赛道的功率、速度、高程与距离对应关系亦以图形方式给出。

\section{模型敏感性分析}

\subsection{对天气与环境的敏感性}

在第 5.4 节中，针对天气与环境因素的影响，我们使用外部因素分析模型将这些影响因素归纳为：气温（$x_1$）、气压（$x_2$）、海拔（$x_3$）、风阻（$x_4$），最终目标是使完成比赛的时间 $t$ 最小。

基于用子集模拟重要性采样的思想，我们将敏感性转化为条件失效概率对底层变量的偏导数。该方法具有较高的计算效率与精度。随着数据集中样本数量的增加，失效概率将逐渐趋近于一个稳定值。

各变量的计算结果如下。对气温 $x_1$，偏导数 $\partial t/\partial x_1$ 的精确值为 $-1.8921\times 10^{-7}$，估计值为 $-1.8886\times 10^{-7}$，变异系数（COV）为 0.0187，误差率为 0.18\%。对气压 $x_2$，$\partial t/\partial x_2$ 的精确值为 $-1.3748\times 10^{-7}$，估计值为 $-1.4294\times 10^{-7}$，变异系数为 0.248，误差率为 3.97\%。对海拔 $x_3$，$\partial t/\partial x_3$ 的精确值为 $-0.002215$，估计值为 $-0.2219$，变异系数为 0.010，误差率为 0.18\%。对风阻 $x_4$，$\partial t/\partial x_4$ 的精确值为 $-6.1065\times 10^{-7}$，估计值为 $-5.8486\times 10^{-7}$，变异系数为 0.046，误差率为 4.22\%。第 5 个变量 $x_5$ 的偏导数精确值为 $1.4345\times 10^{-7}$，估计值为 $1.4687\times 10^{-7}$，变异系数为 0.015，误差率为 2.38\%。

Cov 是统计学中常用的估计值变异系数，反映估计值的收敛性。Error(\%) 是估计值的失效概率。误差率越大，相应 $x_i$ 的敏感性越高。

由敏感性分析的结果可以得出：以 1 为分界线，气压与风阻属于敏感因素，而温度与海拔属于非敏感因素。因此，气压与风阻是对我们的模型影响相对更大的指标。

\subsection{对偏离目标功率的敏感性}

在实际比赛中，尽管骑手会尽力执行模型所制定的功率分配方案，但累积的体力消耗与能量消耗可能影响骑手对赛道上情形的判断，从而扰乱骑手在不同路段上的功率分配。

同一名运动员采用四种策略：功率总体规划模型所推荐的策略、快速起步、缓慢起步与不断变化。我们以三个不同路段的时间作为目标函数计算敏感性，其中 $t_1$ 对应坡道、$t_2$ 对应平路、$t_3$ 对应弯道。

各策略的偏离目标失效概率（\%）如下：推荐策略在坡道、平路、弯道上分别为 0.42、0.23、0.33；快速起步分别为 3.31、1.81、2.60；缓慢起步分别为 1.54、0.84、1.21；不断变化分别为 5.57、3.05、4.38。

因此，为提升骑手表现，我们对功率总体分配模型作了补充：对四种策略进行优先排序，并对不同路段执行策略的重要性进行排序，即推荐策略 $>$ 缓慢起步 $>$ 快速起步 $>$ 不断变化，坡道 $>$ 弯道 $>$ 平路。

\section{模型向团队计时赛的扩展}

在个人计时赛（TT）中，骑手往往以功能性阈值功率（Functional Threshold Power）或其附近的稳定配速为目标。但在团队计时赛（TTT）中，每名骑手都会以明显更高的功率短时间领骑集团。其目标是保持车队速度一致，但高于任何单个骑手独自骑行所能达到的速度。

最佳策略是：首先，六名骑手排成一条直线，车队的骑行方向与风向平行，这样可以确保处于队首的骑手承担大部分风阻。其次，在赛道前 $a\%$ 的路段，由两名爆发力最强、耐力不足的骑手分别领骑坡道与弯道，其余骑手在该路段不受风阻影响。在赛道的 $a\%$ 至 $b\%$ 段，因承担整个车队风阻而体力严重削弱的这两名选手可以慢慢落到车队后方，其余四名体力更充沛的选手轮流处于队首以抵抗风阻。每名骑手领骑 1\,min 后便退到车队后方，但仍须紧跟车队，这样每名骑手每轮约有 3\,min 的低功率输出休息时间。

此外，整个车队的速度应保持一致。领骑者的输出功率 $P_{leader}$ 应等于多少？显然，$P_{leader}$ 必须大于个人计时赛中各路段的输出功率 $P_{ITT}$。因此我们定义 $P_{leader}=c\times P_{ITT}$。为体现不同队员之间的差异，我们首先使用骑手能力评价模型，通过曲线拟合计算每名骑手的能力指标。

\begin{enumerate}
\item[i)] 如果骑手在某路段上不处于领骑位置，我们的模型需要作相应修改：在受力分析图中作出改动（去除空气阻力 AERODYNAMIC DRAG 的影响），并忽略式（9）。
\item[ii)] 在其他情形下，直接使用未经修改的功率总体规划模型。
\end{enumerate}

这样便得到一个以 $a$、$b$、$c$ 为变量的多元非线性规划问题。我们使用相同的算法求解，并计算平均值，得到 $a$ 约为 45\%，$b$ 约为 85\%，$c$ 约为 130\%。

因此，我们可以得到六人车队的团队策略：i) 在赛道前 45\% 的路段，让两名爆发力强而综合实力相对较弱的骑手轮流领骑，每人领骑 5\,min，其输出功率为 ITT 模型的 130\%；ii) 在赛道 45\% 至 85\% 的路段，起初领骑的两名骑手逐渐选择落到后方，其余四名耐力更强的骑手各以 130\% 的输出功率轮流领骑 1\,min；iii) 在赛道最后 15\%，综合能力（如耐力）最强的骑手始终处于领骑位置，直至比赛结束。

\section{模型的优缺点}

\subsection{优点}

\begin{itemize}
\item \textbf{创新性。}我们的模型使用 6 个独立指标来衡量骑手能力，很好地刻画了骑手之间的差异，使模型相对可靠且全面。所选指标与相关领域的最新研究相一致。
\item \textbf{较高的实际应用价值。}为解决不同问题，我们嵌套使用不同的模型，使对每个问题的分析更具说服力。
\item \textbf{直观的可视化。}我们的模型用大量图形展示功率分配方法，直观清晰。
\end{itemize}

\subsection{缺点}

\begin{itemize}
\item \textbf{部分数据可能不够准确。}我们作了较多假设，其中一些假设较难实现。
\item \textbf{算法变量过多。}尽管这样能够全面地考虑问题，但也导致时间复杂度较高。
\end{itemize}

\section*{参考文献}

\begin{enumerate}
\item[{[1]}] http://cyclingtips.com
\item[{[2]}] https://olympics.com/en/news/uci-road-world-championships-2021-cycling-flanders-preview
\item[{[3]}] Leo P, Spragg J, Podlogar T, Lawley JS, Mujika I. Power profiling and the power-duration relationship in cycling: a narrative review. Eur J Appl Physiol. 2022 Feb;122(2):301-316. doi: 10.1007/s00421-021-04833-y. Epub 2021 Oct 27. PMID: 34708276; PMCID: PMC8783871.
\item[{[4]}] Burnley M, Jones AM (2018) Power-duration relationship: physiology, fatigue, and the limits of human performance. Eur J Sports Sci 18(1):112
\item[{[5]}] Barker T, Poole DC, Noble ML, Barstow TJ (2006) Human critical power-oxygen uptake relationship at different pedalling frequencies. Exper Physiol 91(3):621-632
\item[{[6]}] Bassett DR et al (1999) Comparing cycling world hour records, 1967-1996: modeling with empirical data. Med Sci Sports Exerc 31(11):1655-1676
\item[{[7]}] Billat VL et al (2003) The concept of maximal lactate steady state: a bridge between biochemistry, physiology and sport science. Sports Med 33(6):407-426
\item[{[8]}] Borszcz FK et al (2018) Functional threshold power in cyclists: validity of the concept and physiological responses. Int J Sports Med 39(10):737-742
\end{enumerate}

\section*{附录：赛道指导海报文字内容}

\subsection*{骑手能力}

如何评价骑手的能力？我们的答案是绘制该骑手给定的功率—持续时间曲线。他需要以给定的功率输出，而我们测量骑手在该功率下的持续时间，这就得到一个样本点。显然，绘制这条曲线需要大量样本点，也就需要进行大量实验，每次实验采用不同的给定功率并记录持续时间。这在短时间内并不现实，因此我们提供一种带有一定误差但更简便的测试方法：用三次测试分别计算骑手的最大瞬时输出功率（$P_{\max}$）、临界功率（$CP$）与无氧能力（$W'$），再代入下式即可得到同一条曲线。
\[
t=\frac{W'}{P-CP}+\frac{W'}{CP-P_{\max}}
\]
随着测试次数的增加，这条曲线会被不断修正。

\subsection*{指标的测量}

\begin{itemize}
\item 最大瞬时功率（$P_{\max}$）通过 5\,s 最大功率输出测试获得，取 5\,s 的平均值作为计算值，它反映骑手的爆发力。
\item 临界功率（$CP$）通过 3\,min 全力输出测试获得，取最后 30\,s 记录到的骑手平均功率作为计算值，它反映骑手的摄氧能力与骑手的综合计时赛能力。
\item 无氧能量储备（$W'$）通过 5\,min 最大平均功率测试获得。我们记录该值 $\bar{P}$，则 $W'=(\bar{P}-CP)\times 300\,\mathrm{s}$，它反映骑手的耐力。
\end{itemize}

\subsection*{能力指数}

得到该曲线后，我们可以更细致地评价骑手的能力指数。我们选取 6 个量作为骑手的能力指数，其中包括前述三个量：最大瞬时功率（$P_{\max}$）、临界功率（$CP$）、无氧能量储备（$W'$），以及最基本的身高（$H$）、体重（$m$）与控车能力（$BC$）。控车能力 $BC$ 为身高（cm）除以 2.5 倍体重（kg），它表示下坡最高速度与转弯速度之比。得到六项能力指数后，绘制六边形能力图，六边形面积的大小可以反映骑手的综合实力。

\subsection*{骑手的分类}

我们依据能力指数将所有骑手划分为三个方向：1. 冲刺型骑手：$P_{\max}$ 大、爆发力强，擅长平顺路面。2. 爬坡型骑手：$W'$ 大、耐力强，擅长山路骑行，尤其是在爬坡较多的赛段。3. 全能型选手：$CP$ 很高，$P_{\max}$ 不如冲刺型骑手但依然较高，$W'$ 不如爬坡型骑手，擅长个人计时赛，并在团队计时赛中担任主力。这三类骑手都应具有较高的身高与较小的体重，以提升控车能力，从而在下坡或转弯时既能保持高速又能避免摔车。

\subsection*{案例：ALICE}

例如，女骑手 Alice 的六项能力指标如下。可以看出，Alice 适合参加山地更陡的个人计时赛，并适合在团队计时赛中担任主力。

\subsection*{分析与策略}

我们需要依据不同的影响因素制定不同的策略。首先是依据坡度与急弯对赛道进行分段。我们将骑手能力参数以及赛道、环境与天气参数作为输入，模型即可输出各路段最合理的输出功率。我们模型的定性结果是：赛道坡度越大，输出功率越高；在弯道处应提高输出功率并降低速度。在此之前，安全是前提。输出功率分配的优先次序为：上坡 $>$ 弯道 $>$ 平路 $>$ 下坡。骑手应尽可能优先执行上坡与弯道策略。实际执行与计划的偏差 $<10\%$。

环境因素：经过分析，风阻对比赛成绩的影响最为显著，因此我们建议装备团队减小骑行装备的迎风面积，保持骑手的骑行姿态不变，并尽可能降低风阻系数。地面湿度也是团队需要重点关注的要素，因为它直接影响地面滚动阻力系数，进而影响骑手的安全。建议在潮湿赛道上使用厚度较大、摩擦系数较大的轮胎，以在确保安全的同时降低滚动阻力系数。

\subsection*{示例：2021 年日本东京奥运会计时赛赛道}

本次比赛具有平均坡度大、急弯多的特点。如果 Alice 参加该比赛，她将获得很大优势。我们考察比赛当日风力、风向、气压、气温、地面湿度与空气密度，计算风阻与地面滚动阻力，代入模型后输出该计时赛各路段的计划输出功率。Alice 首先需要注意的是，在陡上坡与急弯处必须尽可能按照计划执行。

\subsection*{团队计时赛}

对于团队计时赛，我们建议由两名最具爆发力的冲刺手在比赛前 45\% 轮流领骑，他们可以为整个车队抵抗空气阻力，其输出功率为个人计时赛的 130\%。在赛道 45\% 至 85\% 段，这两名冲刺手消耗了大量体力并逐渐落到后方，其余四名爬坡型与全能型骑手循环轮流领骑。更好的做法是，安排爬坡型骑手在上坡路段领骑、全能型骑手在平路路段领骑。在最后 15\%，应由两名全能型骑手轮流领骑。两名爬坡型骑手需要紧跟车队，最后四名选手几乎同时冲过终点线。需要注意的是，前述赛道、环境与天气因素在团队计时赛中仍应考虑。骑手能力、战术安排、功率分配策略与环境应被全方位地综合考虑。

（团队 2210307）

\end{document}
